% TOP_PROBLEM: {"id": 88, "title": "Hilbert's 14th Problem: Finite Generation of Rings", "category_id": 4, "category": {"id": 4, "name": "algebra", "display_name": "Algebra", "description": "Group theory, ring theory, field theory, and algebraic structures.", "slug": "algebra", "order_index": 4, "created_at": "2026-07-31T15:26:25.671Z"}, "status": "open"}
% FIRST_POSTED: null
% ENTRY_KIND: statement_only
% Imported from: batch02.tex; UnsolvedMath ID 88
\documentclass[11pt]{article}
\usepackage[margin=25mm]{geometry}
\usepackage{fontspec}
\setmainfont{Latin Modern Roman}
\usepackage{amsmath,amssymb,mathtools}
\usepackage{unicode-math}
\setmathfont{Latin Modern Math}
\usepackage{xurl,hyperref}
\hypersetup{colorlinks=true,urlcolor=blue}
\setcounter{secnumdepth}{0}
\setlength{\parindent}{0pt}
\setlength{\parskip}{5pt}
\setlength{\emergencystretch}{3em}
\newcommand{\sref}[2]{\hyperlink{source:#1:#2}{[S#2]}}
\newcommand{\cataloguescope}[1]{\paragraph{Recorded target.}#1}
\begin{document}
% UnsolvedMath ID 88; source catalogue code HIL-014
\setcounter{section}{87}
\section{Hilbert's 14th Problem: Finite Generation of Rings}\label{88}
{\small\sffamily UnsolvedMath ID 88 · Algebra}
\subsection{Definitions and mathematical statement}

\paragraph{Shared notation.}$\mathbb N=\{1,2,\ldots\}$, and $\mathbb Z,\mathbb Q,\mathbb R,\mathbb C$ denote the integers, rationals, reals and complex numbers. Any other conventions are specified locally.


Let $k$ be a field, $V$ a finite-dimensional $k$-vector space, and $G$ a linear algebraic group over $k$ with a rational linear representation on $V$. Write $k[V]=\operatorname{Sym}_k(V^*)$ for the ring of polynomial functions on $V$, identified with $k[x_1,\ldots,x_n]$ after choosing a basis. Let $k[G]$ be the coordinate ring of $G$ and let $\rho^*:k[V]\to k[G]\otimes_k k[V]$ be the pullback on polynomial functions of the action morphism $\rho:G\times V\to V$. Its invariant ring is
\[
k[V]^G=\{f\in k[V]:\rho^*(f)=1\otimes f\}.
\]
This algebraic-group invariance condition expresses $f(gv)=f(v)$ as an identity of regular functions and distinguishes it from invariance under only the possibly finite group $G(k)$. A $k$-algebra $R$ is \emph{finitely generated} if $R=k[f_1,\ldots,f_m]$ for finitely many $f_i\in R$.
\paragraph{Statement recorded in the supplied source.}
For every field $k$, every finite-dimensional $k$-vector space $V$, and every rational linear action of a linear algebraic group $G$ on $V$, is the invariant ring $k[V]^G$ a finitely generated $k$-algebra? \sref{88}{1}\sref{88}{2}
The question recorded here is the linear-representation formulation. It has a negative answer: Nagata gave counterexamples, and Totaro's Theorem 0.1 supplies linear-representation counterexamples over every field. \sref{88}{2}
\subsection{Short English statement}
Is the ring of invariant polynomials for every linear algebraic-group representation generated by finitely many polynomials? The historical universal claim is false.
\subsection{Sources}
\begin{description}
\item[\hypertarget{source:88:1}{S1}] ULAM.AI and the UnsolvedMath Contributors, \emph{UnsolvedMath: A Curated Collection of Open Mathematics Problems}, record 88 (HIL-014). CC BY 4.0. \url{https://huggingface.co/datasets/ulamai/UnsolvedMath}.
\item[\hypertarget{source:88:2}{S2}] Burt Totaro, \emph{Hilbert's fourteenth problem over finite fields, and a conjecture on the cone of curves}, arXiv:0808.0695v1, 5 August 2008, introduction and Theorem 0.1, p. 1; section 1, p. 2, distinguishing the linear-representation question from broader variants. \url{https://arxiv.org/pdf/0808.0695v1}.
\end{description}
\label{end:88}
\end{document}
